\documentclass[12pt,a4paper]{article}

% ---------- Core packages for NCA submission ----------
\usepackage[utf8]{inputenc}
\usepackage[margin=2.5cm]{geometry}
\usepackage{amsmath,amssymb,amsfonts,amsthm}
\usepackage{algorithm}
\usepackage{algpseudocode}
\usepackage{array}
\usepackage{xcolor}
\usepackage{url}
\usepackage{graphicx}
\usepackage{booktabs}
\usepackage{multirow}
\usepackage{tikz}
\usetikzlibrary{positioning,calc,fit,shapes,backgrounds,arrows.meta}
\usepackage{longtable}
\usepackage{pgfplots}
\pgfplotsset{compat=1.15}
\usepackage{subcaption}
\usepackage{float}
\usepackage{bm}
\usepackage[numbers,sort&compress]{natbib}
\usepackage[hidelinks]{hyperref}
\usepackage[capitalize,noabbrev]{cleveref}
\usepackage{placeins}
\usepackage[final]{microtype}

% ---------- Revision color ----------
\definecolor{revcolor}{RGB}{0,0,180}
\newcommand{\rev}[1]{\textcolor{revcolor}{#1}}

% ---------- Line spacing ----------
\usepackage{setspace}
\onehalfspacing

% ---------- Theorem environments ----------
\theoremstyle{definition}
\newtheorem{theorem}{Theorem}[section]
\newtheorem{lemma}[theorem]{Lemma}
\newtheorem{corollary}[theorem]{Corollary}
\newtheorem{definition}[theorem]{Definition}
\newtheorem{assumption}[theorem]{Assumption}
\newtheorem{proposition}[theorem]{Proposition}
\newtheorem{remark}[theorem]{Remark}

% ---------- Mathematical operators ----------
\DeclareMathOperator{\ECE}{ECE}
\DeclareMathOperator{\KL}{KL}
\DeclareMathOperator*{\argmax}{arg\,max}
\DeclareMathOperator{\softmax}{softmax}
\newcommand{\E}{\mathbb{E}}
\newcommand{\R}{\mathbb{R}}
\newcommand{\D}{\mathcal{D}}
\newcommand{\A}{\mathcal{A}}
\newcommand{\B}{\mathcal{B}}
\newcommand{\N}{\mathcal{N}}
\newcommand{\S}{\mathcal{S}}
\newcommand{\G}{\mathcal{G}}
\newcommand{\F}{\mathcal{F}}

\begin{document}

\title{\textbf{Supplementary Material:\\Byzantine-Resilient Stochastic Games for Federated Multi-Cloud Intrusion Detection}}

\author{Roger Nick Anaedevha \and Alexander G. Trofimov \and Yuri V. Borodachev\\
\small National Research Nuclear University \\ 
\small Moscow Engineering Physics Institute (MEPhI), \\ 
\small Moscow, 115409, Russia}

\date{}
\maketitle

\begin{abstract}
This supplementary material complements the main paper by providing comprehensive details, extended analyses, and complete results. It includes an expanded survey of related work with additional baseline methods and recent developments; a detailed privacy analysis—covering regulatory compliance mapping and differential-privacy budget allocation strategies; and complete mathematical proofs for all theorems establishing existence, uniqueness, privacy, robustness, and convergence guarantees. We further provide extended adversarial-training procedures with domain-specific perturbation strategies; full experimental configurations to ensure reproducibility and fair baseline comparison; additional experimental results (precision–recall curves, ROC analysis, and detailed confusion matrices); scalability analyses with computational-complexity breakdowns; and deployment considerations, including cost–benefit analysis and integration guidelines. Finally, we present a comprehensive limitations analysis with failure-mode discussions and mitigation strategies, and we outline several directions for future research.
\end{abstract}


\tableofcontents
\clearpage

\section{Appendix A: Extended Related Work}

\subsection{A.1 Extended Game-Theoretic Security Literature}

Beyond the foundational works discussed in the main paper, several recent developments inform our approach to game-theoretic network security. Roy et al. extended classical two-player security games to multi-player scenarios, introducing security games with multiple defenders. Their analysis revealed the complexity of coordination in distributed defense, showing that Nash equilibria may be inefficient due to lack of cooperation. The work established that the price of anarchy in security games can reach order n where n is the number of defenders, highlighting the need for coordination mechanisms in federated settings. Our framework addresses this challenge through explicit coordination mechanisms built into the federated aggregation protocol in Algorithm 1 of the main paper.

The work of Alpcan and Başar on stochastic security games provides the mathematical foundation for modeling dynamic interactions under uncertainty. They introduced saddle-point equilibria in security games, demonstrating that optimal defense strategies exist under certain regularity conditions. However, their approach assumes complete information and centralized decision-making, limitations that become critical in distributed cloud environments where each defender has only local visibility.

Recent advances by Zhu and Başar incorporate dynamic aspects through differential games, establishing Hamilton-Jacobi-Bellman equations for optimal defense strategies. Their framework models continuous-time interactions but requires complete state observability, which is unrealistic in multi-cloud environments. Moreover, their computational approach scales exponentially with state space dimension, limiting practical applications. Our jump-diffusion formulation in Equations 6 through 8 of the main paper provides a more scalable alternative.

\subsection{A.2 Advanced Federated Learning for Security}

Preuveneers et al. first applied federated learning to intrusion detection, achieving 97.8 percent accuracy on the NSL-KDD dataset. Their approach demonstrated the feasibility of distributed security model training but did not consider adversarial participants, privacy guarantees, or the heterogeneous characteristics of modern cloud environments.

Wang et al. introduced CloudFL, a federated learning framework specifically designed for multi-cloud security, achieving 92.1 percent accuracy. However, CloudFL does not incorporate game-theoretic considerations to model strategic adversarial behavior and assumes relatively mild heterogeneity compared to the extreme diversity in ICS3D datasets.

Chen et al. developed AdaptiveFL-Container with 89.7 percent detection rates for container security, introducing domain adaptation techniques. Their approach focuses specifically on container security and does not generalize to other cloud domains. Liu et al. proposed hierarchical federated learning for IoT security, demonstrating 91.3 percent accuracy. Their hierarchical architecture addresses scalability concerns but does not provide Byzantine resilience guarantees.

\subsection{A.3 Extended Byzantine Robustness Analysis}

The vulnerability of federated learning to Byzantine attacks was first systematically studied by Blanchard et al., who showed that a single malicious client can arbitrarily manipulate the global model in standard FedAvg implementations. Their Krum algorithm provides some robustness but fails under coordinated attacks.

Yin et al. introduced coordinate-wise median and trimmed mean aggregation, achieving breakdown points near 50 percent. However, their analysis assumes homogeneous data distributions and may fail when natural heterogeneity across domains mimics Byzantine behavior.

Recent advances include the BALANCE algorithm by Zhang et al., which enables each client to use its local model as a similarity reference. BALANCE achieves strong theoretical guarantees under IID assumptions but degrades with extreme class imbalance common in security applications.

Zhou et al. introduced RobustFL with 94.2 percent accuracy under 30 percent Byzantine corruption. Patel et al. developed SecureAgg+ achieving 93.8 percent robustness with improved communication efficiency. Singh et al. proposed Byzantine-resilient federated learning for edge computing with 92.6 percent performance retention.

\section{Appendix B: Privacy Analysis and Regulatory Compliance}

\subsection{B.1 Extended Differential Privacy Analysis}

\begingroup
\setlength{\emergencystretch}{2em}

Our framework achieves \((\epsilon_d,\delta_d)\) differential privacy per domain through the composition of gradient clipping and the Gaussian mechanism. Consider an adaptive adversary that observes all previous federated rounds and selects Byzantine clients accordingly. Under our domain-adaptive privacy mechanism, the cumulative privacy loss after \(T\) rounds satisfies
\[
\epsilon_{\text{total}}
\;\le\;
\epsilon_d\,\sqrt{\,2T\ln(1/\delta_d)\,}\;+\;T\,\delta_d.
\]

The proof follows from advanced composition theorems for differential privacy. Each federated round adds independent Gaussian noise calibrated to privacy budget \(\epsilon_d\). Under adaptive composition, the total privacy loss grows as \(O(\sqrt{T})\) rather than linearly, due to the Gaussian mechanism’s strong composition properties.

The sensitivity analysis requires careful treatment of domain-specific characteristics. The clipping operation bounds the \(\ell_2\) sensitivity as
\[
\Delta_{2}\;\le\;2\,C_d\,\rho_d^{-1/2}.
\]
The imbalance-adjusted clipping norm accounts for varying gradient magnitudes across domains with different imbalance ratios.

\endgroup

\subsection{B.2 Regulatory Compliance Mapping}

\begingroup
\setlength{\emergencystretch}{2em}

We provide detailed mapping of our privacy mechanisms to specific regulatory requirements. For GDPR compliance, Article~5 paragraph~1 subparagraph~(f) requires integrity and confidentiality in data processing. Our differential privacy mechanism ensures that individual client data cannot be reverse-engineered from aggregate model updates. Article~25 requires data protection by design, which our domain-specific privacy budgets satisfy by building privacy protections into the system architecture from the outset.

For HIPAA compliance relevant to SOC data containing health entities, Section~164.312 requires transmission security and access control. Secure aggregation with cryptographic commitments ensures transmission security. Entity embedding layers obfuscate individual identifiers, preventing unauthorized access.

For SOX compliance applicable to financial services, Section~404 requires assessment of internal controls. Byzantine detection and reputation tracking provide auditable internal controls over model training with complete traceability for compliance audits.

\endgroup

\subsection{B.3 Privacy Budget Allocation Strategies}

\begingroup
\setlength{\emergencystretch}{2em}

Different deployment scenarios require different privacy budget allocation strategies. Uniform allocation divides total privacy budget \(\epsilon_{\text{total}}\) uniformly across \(T\) rounds as
\[
\epsilon_{\text{round}} \;=\; \frac{\epsilon_{\text{total}}}{\sqrt{T}}.
\]
This approach is simple but may be suboptimal when attack distributions vary over time.

Adaptive allocation assigns larger budgets to rounds with higher attack prevalence detected through temporal analysis. Domain-prioritized allocation assigns budgets proportionally to domain sensitivity as
\[
\epsilon_{\text{SOC}} \;<\; \epsilon_{\text{Container}} \;<\; \epsilon_{\text{Edge}},
\]
reflecting varying privacy requirements. Our experiments use domain-prioritized allocation with
\(\epsilon_{\text{SOC}}=1.8\), \(\epsilon_{\text{Container}}=2.0\), and \(\epsilon_{\text{Edge}}=2.5\).

\endgroup

\section{Appendix C: Complete Mathematical Proofs}

\subsection{C.1 Proof of Theorem 1: Existence and Uniqueness}

\begingroup
\setlength{\emergencystretch}{2em}

We establish existence and uniqueness of solutions to the coupled system of stochastic differential equations through Picard iteration in the space of càdlàg processes. Assume the drift coefficients satisfy the Lipschitz condition
\[
\big\|\mu_d(s,a,t)-\mu_d(s',a,t)\big\| \;\le\; L_{\mu}\,\|s-s'\|,
\]
and the diffusion coefficients satisfy a Lipschitz condition with constant \(L_{\sigma}\). Assume the compensator measures satisfy integrability conditions
\[
\int z^{2}\,\nu_{i}^{(d)}(dz)\;<\;\infty \quad \text{for all } i,d.
\]
Define the Picard iteration sequence beginning with \(s_{t}^{(d,0)}=s_{0}^{(d)}\), then recursively defining subsequent terms.

Under the Lipschitz conditions, the expected value of the supremum norm squared of successive differences is bounded by \(C\) times the integral of the expected norm squared of previous differences. By Grönwall’s inequality, this establishes convergence of the Picard sequence.

For uniqueness, suppose \(s_t\) and \(s'_t\) are two solutions. Define \(\Delta_t=s_t-s'_t\). Taking expectations and applying the Lipschitz conditions yields
\[
\mathbb{E}\big[\|\Delta_t\|^{2}\big] \;\le\; C\int_{0}^{t}\mathbb{E}\big[\|\Delta_u\|^{2}\big]\,du.
\]
Grönwall’s inequality gives \(\mathbb{E}\big[\|\Delta_t\|^{2}\big]=0\), implying uniqueness almost surely.

\endgroup

\subsection{C.2 Proof of Theorem 2: Nash Equilibrium Existence}

\begingroup
\setlength{\emergencystretch}{2em}

We apply Kakutani’s fixed-point theorem to the best-response correspondence. Define mixed strategy spaces \(\Delta_k^{(d)}\) as probability distributions over action spaces; these are compact and convex subsets of finite-dimensional Euclidean space.

The imbalance-weighted payoff can be written as
\[
u_k^{(d)} \;=\; \sum_{c} w_c^{(d)} \sum_{(s,a)} \pi(\,s,a\,)\,R_k(s,a),
\]
and the best-response correspondence as
\[
\mathrm{BR}_k^{(d)} \;=\; \argmax_{\pi_k^{(d)}} \; u_k^{(d)}.
\]
This correspondence is non-empty (a maximum of a continuous function over a compact set is attained), convex-valued (payoff is linear in \(\pi_k^{(d)}\)), and upper hemicontinuous by the maximum theorem. The product correspondence \(\mathrm{BR}\) maps the compact convex set to itself; by Kakutani’s fixed-point theorem, there exists a fixed point \(\pi^{\star}\), which is a Nash equilibrium. To establish the imbalance-robustness bound, note that imbalance weights introduce regularization; the Kullback–Leibler divergence between the equilibrium and the uniform distribution is bounded, and by Pinsker’s inequality this yields the stated bound.

\endgroup

\subsection{C.3 Proof of Theorem 3: Privacy and Robustness Guarantees}

\begingroup
\setlength{\emergencystretch}{2em}

We prove privacy and robustness separately. For privacy, the mechanism consists of gradient clipping followed by Gaussian noise addition. The clipping bounds \(\ell_2\) sensitivity as
\[
\Delta_{2}\;\le\;2\,C_d\,\rho_d^{-1/2}.
\]
Adding Gaussian noise achieves \((\epsilon_d,\delta_d)\) differential privacy when \(\sigma_d\) is appropriately calibrated. When aggregating over \(n_d\) honest clients, the effective noise scale becomes \(\sigma_d/\sqrt{n_d}\), yielding the stated privacy bound.

For robustness, the trimmed mean removes an \(\alpha_d\) fraction of extreme values. Under the assumption \(f_d < \alpha_d/2\), all Byzantine contributions are removed. The honest gradients follow
\[
g_k \;=\; \nabla \mathcal{L} + \xi_k,
\]
with zero-mean noise of variance \(O(\rho_d)\). The trimmed mean satisfies a concentration bound in which
\[
\big\|\overline{g}-\nabla \mathcal{L}\big\|
\;\le\; O\!\big(\alpha_d \sqrt{\rho_d}\big)\;+\;\text{averaging error},
\]
and adding DP noise contributes an additional error term. Combining yields the total robustness bound.

\endgroup

\subsection{C.4 Proof of Theorem 4: Almost-Sure Convergence}

\begingroup
\setlength{\emergencystretch}{2em}

We use martingale convergence techniques with the Lyapunov function from Definition~6 and consider the expected decrease in this Lyapunov function. By the update rule
\[
\theta_{k}^{(d,t+1)}=\theta_{k}^{(d,t)}-\eta_{d}^{(t)}\,g_{k}^{(d,t)},
\]
we expand the squared norm. Under Byzantine resilience, the aggregated gradient satisfies
\[
\mathbb{E}\big[\tilde g^{(t)}\big]
= \nabla \mathcal{L}\!\left(\theta^{(t)}\right)
+ O\!\big(\alpha_{d}\sqrt{\rho_{d}}\big).
\]
Assume the loss is \(L\)-smooth. At Nash equilibrium, the gradient satisfies the equilibrium condition. Combining these yields an expected descent bound. With domain-adaptive learning rates
\[
\eta_{d}^{(t)}=\eta_{0}\,\rho_{d}^{1/2}/(t+1)^{2/3},
\]
the descent term dominates for large \(t\). Define a martingale \(M_{t}\) accordingly. By the martingale convergence theorem, \(M_{t}\) converges almost surely, implying convergence of the parameter sequence.

\par\endgroup


For convergence rate, sum descent inequality over rounds. With learning rate order 1 divided by t to the power two-thirds, sums evaluate to order of T to the power one-third, establishing the stated convergence rate.

\section{Appendix D: Extended Adversarial Training Procedures}

\subsection{D.1 Protocol-Aware Perturbations for Edge-IIoT}

For MQTT protocol perturbations, we perturb topic lengths within valid ranges from 1 to 255 bytes by adding uniform random noise scaled by perturbation budget, then clipping to maintain protocol validity. We modify message types by randomly swapping with adjacent valid types in the set containing CONNECT, PUBLISH, SUBSCRIBE, ensuring perturbed message type remains valid. We adjust Quality of Service levels by adding uniform random values from the set containing negative 1, 0, and 1, then taking modulo 3 to wrap around.

Modbus function codes range from 1 to 127 with specific codes for operations such as read coils at hexadecimal 01, read discrete inputs at hexadecimal 02, write single coil at hexadecimal 05. We perturb function codes by adding uniform random values, then taking modulo 127 ensuring result remains in valid range. For address manipulation, coil and register addresses are 16-bit values. We add Gaussian noise with variance proportional to squared perturbation budget, then clip to valid range.

HTTP protocol perturbations target request method encoding, content-length fields, and header field variations. We manipulate request method encoding by introducing case variations and whitespace insertions that remain technically valid. Content-length fields are perturbed by adding small integer values maintaining consistency with payload sizes. Header field variations include reordering headers and introducing duplicates, all while maintaining HTTP protocol validity.

\subsection{D.2 CVE-Targeted Adversarial Training for Containers}

CVE-2019-5736 represents a critical vulnerability allowing attackers to overwrite the host runc binary via /proc/self/exe symbolic link, achieving container escape. Attack signatures include elevated access patterns to /proc filesystem, unusual file descriptor manipulation, and exec system calls with suspicious arguments. Our adversarial perturbations inject /proc/self/exe access pattern signatures that mimic benign monitoring while including subtle exploitation anomalies. We modify flow duration to match escape timing profiles ranging from 100 to 500 milliseconds for critical exploitation phase. We adjust packet sizes to reflect privilege escalation attempts producing packets in range 64 to 512 bytes.

CVE-2021-25741 is a path traversal vulnerability in Kubernetes allowing attackers to access files outside volume mounts. Attack indicators include kubelet API call patterns with unusual path parameters, volume mount access sequences attempting filesystem escape, and symlink traversal indicators. Our adversarial training simulates kubelet API call patterns mimicking benign volume operations while including edge cases in path handling. We modify bidirectional flow statistics to match patterns during path traversal attempts showing asymmetric request-response patterns.

Node-RED RCE vulnerabilities enable arbitrary code execution through dashboard exploits and API vulnerabilities. Exploitation patterns include dashboard access pattern variations, HTTP payload sizes matching attack signatures between 256 and 2048 bytes, and flow inter-arrival times adjusted to evade detection. Our adversarial examples generate dashboard access patterns exploring the boundary between legitimate administration and malicious exploitation.

\subsection{D.3 Entity-Aware Perturbations for SOC Domain}

SOC data contains complex entity relationships spanning 33 distinct types including DeviceId, AccountSid, IpAddress. Adversarial perturbations must maintain entity type consistency while disrupting detection. For entity masquerading, we extract primary entities and generate perturbations by selecting similar entities using cosine similarity computed over entity embedding spaces. We select replacement entities with similarity exceeding 0.7 to ensure perturbed incidents remain plausible while evading detection.

Temporal correlation disruption targets multi-window aggregation over 1-hour, 24-hour, and 7-day windows. We shift event timestamps by random amounts within window boundaries to break correlation detection. We add noise to window-based features such as alert counts and average severity, with noise magnitude proportional to window size. We modify temporal features to match normal operational rhythms despite malicious intent, ensuring attack events align with typical business hours.

Alert pattern camouflage modifies alert characteristics to blend with benign activity. We adjust alert severity within reasonable bounds toward borderline values less likely to trigger immediate investigation. Suspicion scores are perturbed to values just below typical detection thresholds. Threat family classifications are modified subtly by introducing confusion between related threat categories.

\section{Appendix E: Detailed Experimental Configurations}

\subsection{E.1 Baseline Configuration and Fairness Protocol}

To ensure fair comparison, all baseline methods underwent identical hyperparameter tuning using the same computational budget. The three-stage tuning protocol begins with coarse grid search allocating 16 GPU-hours to explore learning rates in the set containing 0.0001, 0.0005, 0.001, 0.005, and 0.01, batch sizes in the set containing 64, 128, 256, 512, and 1024, and local epochs in the set containing 1, 3, 5, and 10. Validation on 10 percent held-out data determines the best configuration per method.

Stage two performs fine-grained refinement allocating 16 GPU-hours to narrow search around Stage 1 optimum. Learning rate is refined in plus or minus 50 percent range with 5 logarithmically-spaced values. Momentum parameters and weight decay are tuned independently. Method-specific parameters such as FedProx proximal term mu, SCAFFOLD control variates, and BALANCE similarity threshold are tuned using method-appropriate ranges.

Stage three performs 5-fold cross-validation allocating 16 GPU-hours on each domain with best Stage 2 configuration. Final hyperparameters are selected based on mean validation accuracy across domains weighted by domain size. This ensures selected hyperparameters generalize rather than overfitting to particular validation splits.

\begingroup
\setlength{\emergencystretch}{2em}

The optimal hyperparameters show most methods perform best with base learning rate \(0.001\), batch size \(256\), and \(5\) local epochs. All methods use Adam optimizer with momentum \(\beta_{1}=0.9\) and \(\beta_{2}=0.999\), weight decay \(0.001\), and gradient clipping norm \(1.0\). Method-specific parameters include FedProx \(\mu=0.01\), SCAFFOLD control variate learning rate \(0.1\), BALANCE similarity threshold \(0.7\), CloudFL domain adaptation weight \(0.3\), and RobustFL trim ratio \(0.1\).

\par\endgroup


\subsection{E.2 Hyperparameter Selection Justification}

\begingroup
\setlength{\emergencystretch}{2em}

The domain-adaptive learning rate formula \(\eta_d^{(t)}=\eta_0\,\rho_d^{1/2}/(t+1)^{2/3}\) emerged from both theoretical analysis and empirical validation. From convergence analysis in Theorem 4, the optimal learning rate decay for non-convex objectives with heterogeneous data follows order of \(1/t^{\alpha}\) where \(\alpha\) is between \(1/2\) and \(2/3\). The two-thirds power provides best balance between exploration and exploitation.

The imbalance scaling through \(\rho_d^{1/2}\) scales learning rates proportionally to domain imbalance severity. Domains with higher imbalance require larger learning rates to ensure minority classes contribute meaningfully despite having fewer training examples. This scaling emerged from grid search over \(\rho_d^{\alpha}\) for \(\alpha \in \{0.25,0.5,0.75,1.0\}\), with \(\alpha = 0.5\) achieving optimal performance.

\par\endgroup


The Dirichlet concentration parameter alpha equals 0.3 for data partitioning simulates realistic heterogeneity observed in multi-organizational deployments. Preliminary studies comparing alpha values revealed that 0.3 best matches label distribution divergence observed in real-world federated security systems. Jensen-Shannon divergence shows 0.47 plus or minus 0.08 for alpha equals 0.3 compared to 0.42 plus or minus 0.05 observed in production deployments.

Domain-specific privacy budgets were chosen through consultation with industry partners and regulatory compliance experts. SOC environments with epsilon equals 1.8 represent most sensitive domain due to entity relationships, chosen to satisfy European Data Protection Board guidance suggesting epsilon less than 2 for personal data. Container environments with epsilon equals 2.0 reflect moderate sensitivity. Edge-IIoT environments with epsilon equals 2.5 face less stringent requirements due to primarily operational data.

\subsection{E.3 Reproducibility Details}

All experiments use fixed random seeds for complete reproducibility. NumPy seed is set to 42, PyTorch seed to 123, CUDA seed to 456, Python random seed to 789, and TensorFlow seed to 101 when applicable. These seeds control all sources of randomness including weight initialization, data shuffling, dropout operations, and stochastic gradient sampling.

The software environment consists of Python 3.9.7, PyTorch 1.12.1, CUDA 11.6, cuDNN 8.4.0, NumPy 1.23.0, Pandas 1.4.3, Scikit-learn 1.1.1, Matplotlib 3.5.2, and Seaborn 0.11.2. Version pinning ensures consistent behavior across different execution environments.

Hardware specifications include four NVIDIA A100 GPUs each with 40GB VRAM, two AMD EPYC 7742 CPUs each with 64 cores running at 2.25GHz, 512GB DDR4-3200 RAM, 4TB NVMe SSD with PCIe 4.0 interface, and 100 Gbps Infiniband networking for multi-GPU communication.

Statistical testing methodology ensures robustness. Each result represents mean plus or minus standard deviation across 5 independent runs with seeds 42, 123, 234, 345, and 456. Paired t-tests compare methods on identical data splits. Bonferroni correction adjusts significance threshold to alpha equals 0.05 divided by m where m is number of comparisons. Cohen's d effect size quantifies practical significance. All improvements achieve p less than 0.01 after correction with Cohen's d greater than 2.5.

\section{Appendix F: Additional Experimental Results}

\subsection{F.1 Precision-Recall Analysis}

We provide high-resolution precision-recall curves for all three domains demonstrating FedGTD's superior performance across the full operating range from low recall high precision operating points to high recall lower precision operating points appropriate when missing attacks carries severe consequences.

The Edge-IIoT domain with moderate 2.67:1 imbalance shows FedGTD achieving PR-AUC of 0.937 compared to 0.891 for RobustFL and 0.847 for FedAvg, representing 5.2 percent improvement over best baseline. The Container domain with severe 15.7:1 imbalance shows FedGTD achieving PR-AUC of 0.924 compared to 0.878 for RobustFL, representing 5.2 percent improvement. The SOC domain with extreme 99:1 imbalance shows most dramatic improvement with FedGTD achieving PR-AUC of 0.687 compared to 0.531 for RobustFL, representing 29.4 percent relative improvement validating effectiveness of imbalance-weighted objectives under extreme conditions.

\subsection{F.2 Detailed Confusion Matrices}

The Edge-IIoT confusion matrix reveals balanced performance across all 14 attack families despite severe class imbalance. Rare attacks like XSS at 0.8 percent of data achieve 88.7 percent recall, while common attacks like DDoS maintain 96.4 percent recall. The primary confusion occurs between similar attack types, for example SQL Injection and XSS showing 1.2 percent cross-confusion, reflecting semantic similarity in web-based exploit patterns. Normal traffic achieves 97.2 percent precision demonstrating low false positive rate crucial for operational deployment.

The Container confusion matrix demonstrates precise CVE-specific detection with minimal cross-confusion. Container escape CVE-2019-5736 achieves 96.7 percent recall with only 1.1 percent confusion with path traversal attacks, where both involve privilege escalation mechanisms explaining limited overlap. The 2.2 percent false positive rate on benign traffic represents acceptable operational overhead for critical container security.

The SOC confusion matrix addresses the most challenging scenario with 99:1 overall imbalance where true positive incidents represent only 0.8 percent of all incidents. FedGTD achieves 92.3 percent recall on true positives, substantially exceeding RobustFL at 68.8 percent. Critically, 95.6 percent of misclassified true positives are labeled as benign positive rather than false positive, ensuring these incidents still receive analyst attention. This confusion pattern aligns with operational SOC workflows where distinguishing true positive from benign positive requires detailed investigation.

\subsection{F.3 Statistical Significance Testing}

All tests compare FedGTD versus RobustFL as best baseline on identical data splits with 5 independent runs. For Edge-IIoT Accuracy with df equals 4, t-statistic equals 8.34, p-value equals 0.0012, and Cohen's d equals 3.73. For Container Accuracy, t-statistic equals 7.21, p-value equals 0.0019, Cohen's d equals 3.22. For SOC Accuracy, t-statistic equals 12.47, p-value equals 0.0002, Cohen's d equals 5.57.

For Edge Macro-F1, t-statistic equals 9.12, p-value equals 0.0008, Cohen's d equals 4.08. For Container Macro-F1, t-statistic equals 6.89, p-value equals 0.0024, Cohen's d equals 3.08. For SOC Macro-F1, t-statistic equals 11.23, p-value equals 0.0003, Cohen's d equals 5.02.

For Byzantine Retention at 20 percent, t-statistic equals 5.67, p-value equals 0.0047, Cohen's d equals 2.53. For Communication Efficiency, t-statistic equals 14.32, p-value equals 0.0001, Cohen's d equals 6.40. All improvements remain statistically significant after Bonferroni correction for 12 comparisons with threshold alpha equals 0.00417. Cohen's d values ranging from 2.53 to 6.40 demonstrate very large practical significance.

\section{Appendix G: Extended Scalability Analysis}

\subsection{G.1 Computational Complexity Analysis}

\begingroup
\setlength{\emergencystretch}{2em} % extra flexibility; avoids overfull boxes
the domain-adaptive learningcomplexity for each client performing \(E\) local epochs over a dataset of size \(n_k^{(d)}\) with model dimension \(d\) equals order of \(E\,n_k^{(d)}\,d\,L\), where \(L\) is number of layers. Gradient clipping requires order of \(d\) operations. Byzantine detection computing pairwise gradient similarities requires order of \(K^{2}d\) operations, representing primary scalability bottleneck. Trimmed mean aggregation requires order of \(d\,K\log K\) operations.

Total per-round complexity combines all components as
\[
O\!\left(E\,n_{\max}\,d\,L\right)\;+\;O\!\left(K^{2}d\right)\;+\;O\!\left(d\,K\log K\right).
\]
For our experimental configuration with \(K=20\), \(d\approx 10^{6}\), and \(E=5\), the Byzantine detection term dominates aggregation overhead.
\par\endgroup


\subsection{G.2 Scalability to Larger Federations}

We conducted experiments assessing scalability to larger federations with K in the set containing 20, 40, 60, 80, and 100 participants. Detection accuracy improves modestly from 96.1 percent at K equals 20 to 96.6 percent at K equals 100, a gain of 0.5 percentage points due to better attack coverage. Per-round time grows approximately quadratically from 39.7 seconds at K equals 20 to 358.1 seconds at K equals 100, with growth exponent 1.95 closely matching theoretical K squared complexity.

Communication overhead per round grows linearly from 1.92 gigabytes at K equals 20 to 9.60 gigabytes at K equals 100. Byzantine detection time grows quadratically from 15.3 seconds to 467.3 seconds, confirming this component as bottleneck. Memory usage grows approximately linearly from 15.2 gigabytes to 72.3 gigabytes.

For K equals 100 participants, per-round time reaching 358 seconds or approximately 6 minutes may be acceptable for many deployment scenarios where federated rounds occur hourly or daily, but could be improved through approximation techniques for more frequent updates or larger federations.

\subsection{G.3 Proposed Scalability Improvements}

Hierarchical federation architecture addresses order of K squared Byzantine detection complexity through multi-tier organization. Level 1 organizes clients into clusters of size k approximately equal to square root of K. Each cluster performs Byzantine detection locally with complexity order of k squared times d. Level 2 regional coordinators receive cluster aggregates and perform second-level detection with complexity order of K times d. Total complexity becomes order of K times d instead of order of K squared times d, enabling scaling to thousands of participants.

Approximate Byzantine detection uses sampling and hashing techniques. Locality-sensitive hashing projects gradients to lower-dimensional space d-prime much less than d, computing similarities in reduced space with complexity order of K times d-prime. Random sampling has each client compare against s much less than K randomly sampled clients, reducing complexity to order of K times s times d where s can be chosen as order of logarithm of K.

Preliminary experiments with LSH-based approximation using d-prime equals order of logarithm of d show 2 to 3 percentage point degradation in Byzantine detection accuracy but enable scaling to K equals 500 with per-round times under 120 seconds, representing 50 times speedup compared to exact pairwise similarity computation.

\section{Appendix H: Deployment Considerations and Cost Analysis}

\subsection{H.1 Detailed Cost-Benefit Analysis}

We provide comprehensive cost-benefit analysis based on three pilot deployments anonymized as Site A operating primarily SOC, Site B focused on Edge-IIoT, and Site C with multi-domain deployment. Initial deployment costs include model training consuming 1200 dollars at Site A, 800 dollars at Site B, and 1600 dollars at Site C for GPU compute. Integration development ranges from 12000 to 25000 dollars depending on complexity of existing security infrastructure. System testing requires 6000 to 12000 dollars ensuring compatibility. Staff training costs 4000 to 8000 dollars preparing security analysts. Infrastructure upgrades range from 3000 to 8000 dollars. Total initial costs range from 27800 dollars at Site B to 54600 dollars at Site C.

Annual operational costs include federated training GPU compute ranging from 3200 to 6400 dollars depending on domain coverage. Network bandwidth costs 1800 to 3600 dollars annually. System maintenance requires 5000 to 10000 dollars for software updates. Monitoring and support costs 6000 to 12000 dollars for ongoing technical assistance. Total annual operational costs range from 16000 dollars at Site B to 32000 dollars at Site C.

Quantified annual benefits include false positive reduction saving 120000 to 250000 dollars through reduced analyst workload, with 67 percent reduction in false positives directly translating to proportional reduction in investigation time. Faster incident response saves 60000 to 140000 dollars through 23 percent reduction in mean time to respond. Novel attack detection saves 180000 to 450000 dollars through 89 percent improvement in detecting previously unseen attack variants. Compliance automation saves 30000 to 75000 dollars through reduced manual audit effort. Reduced breach risk quantified at 100000 to 250000 dollars annually accounts for decreased probability of costly security incidents. Total annual benefits range from 490000 dollars at Site B to 1165000 dollars at Site C.

Return on investment calculated as annual benefit minus initial cost minus annual cost, all divided by initial cost plus annual cost, yields 1018 percent at Site B, 1380 percent at Site A, and 1248 percent at Site C. All sites achieve positive ROI within first year with break-even occurring at 4 to 5 months.

\subsection{H.2 Integration Guidelines}

The step-by-step integration process begins with assessment phase during weeks 1 through 2, inventorying existing security tools and data sources, identifying data export and API capabilities supporting formats like NetFlow, IPFIX, CEF, and LEEF, assessing network bandwidth and computational resources against requirements, and defining privacy requirements and regulatory constraints applicable to the organization.

Pilot deployment during weeks 3 through 6 deploys FedGTD in shadow mode parallel to existing systems without affecting production security operations, configures domain-specific preprocessing pipelines adapting to organizational data characteristics, establishes baseline performance metrics for comparison, and tunes hyperparameters for organization-specific data distributions.

Production rollout during weeks 7 through 10 transitions to active threat detection with FedGTD alerts feeding into security workflows, integrates with existing SIEM and SOAR platforms through standard APIs, configures alert routing and escalation policies aligned with organizational procedures, and trains security analysts on interpreting new alerts.

Federation expansion during week 11 and beyond identifies partner organizations for federated collaboration, establishes data sharing agreements and privacy protocols satisfying legal and competitive requirements, configures cross-organizational federated learning with appropriate privacy budgets, and monitors cross-domain knowledge transfer effectiveness measuring improvement from collaboration.

\subsection{H.3 Operational Best Practices}

Model update schedules balance timeliness with computational cost. Continuous learning performs federated rounds every 6 to 12 hours for dynamic threat landscapes where attack patterns evolve rapidly. Batch updates with daily federated rounds suit stable environments with predictable attack patterns. Emergency updates trigger immediate federated rounds upon novel attack detection, enabling rapid dissemination of threat intelligence across the federation.

Byzantine client handling establishes procedures for addressing persistent low-reputation clients. Reputation scores reviewed monthly by security team identify clients consistently providing suspicious updates. Automatic quarantine activates for clients with reputation below 0.3, temporarily excluding them from federated learning until investigation completes. Manual investigation required for persistent low-reputation clients determines whether issues stem from data quality problems, misconfiguration, or malicious behavior.

Privacy budget management monitors cumulative privacy loss across all federated rounds, alerting when approaching domain-specific privacy budgets to enable proactive management. Quarterly privacy budget resets after data retention expiration periods leverage legal requirements to delete old data, enabling sustainable long-term operation. Adaptive budget allocation based on attack prevalence directs more budget to high-threat periods while conserving budget during quiet periods.

\section{Appendix I: Comprehensive Limitations and Future Work}

\subsection{I.1 Detailed Limitations Analysis}

Our convergence analysis assumes bounded gradient variance where expected norm squared of gradient difference is bounded. This assumption can fail when industrial IoT devices undergo firmware updates that fundamentally alter traffic patterns. For example, firmware update changing MQTT keepalive intervals from 60 seconds to 300 seconds creates distribution shift violating the variance bound. Sophisticated attackers can intentionally induce concept drift by gradually modifying attack patterns.

Mitigation strategies include implementing online change-point detection to identify concept drift events through statistical tests on gradient distributions, temporarily increasing learning rates and privacy budgets during detected drift periods to enable rapid adaptation, using importance weighting to down-weight clients experiencing drift based on detected distributional changes, and maintaining historical models enabling selective reversion when drift is detected.

We require loss functions to be L-Lipschitz continuous. Neural networks with ReLU activations have discontinuous gradients at zero. While violation is mild for typical weight initializations, adversarial perturbations can exploit these discontinuities. Carefully crafted adversarial inputs can cause large gradient changes with small parameter perturbations, locally violating Lipschitz continuity.

Mitigation strategies include using smooth activations such as GELU or Swish instead of ReLU to ensure differentiability, applying gradient clipping more aggressively near discontinuities detected through gradient magnitude monitoring, implementing adversarial training to smooth the loss landscape, and monitoring gradient magnitude distributions flagging anomalies for investigation.

\begingroup
\setlength{\emergencystretch}{2em}

Our Byzantine resilience guarantees require fewer than \(\alpha_d/2\) malicious clients per domain. When this threshold is exceeded, Byzantine clients can coordinate to bias the trimmed mean. In our experiments with \(40\%\) corruption—exceeding the threshold—performance degrades from \(96.3\%\) to \(71.4\%\). Unlike catastrophic failures in non-robust methods, FedGTD maintains partial functionality through graceful degradation.

\par\endgroup

\subsection{I.2 Generalization Limitations}

Industrial Control Systems impose stringent false positive constraints below 0.1 percent because false alarms can trigger emergency shutdowns costing millions of dollars. Our 3.7 percent false positive rate exceeds this threshold, limiting direct applicability to safety-critical ICS without additional filtering. Required improvements include implementing multi-stage verification with human-in-the-loop for shutdown-triggering alerts, developing ICS-specific thresholds calibrated to operational risk tolerance, integrating with process control models to validate alert plausibility, and using ensemble methods combining FedGTD with domain-specific rule engines.

5G networks involve ultra-low latency requirements below 1 millisecond for Ultra-Reliable Low-Latency Communications services. Our 15 millisecond computational overhead exceeds this requirement. Additionally, network slicing creates dynamic topologies that violate our static domain assumptions. Adaptation requirements include developing hardware-accelerated inference for sub-millisecond latency, designing dynamic domain adaptation for ephemeral network slices, integrating with 5G Management and Orchestration frameworks, and implementing predictive pre-computation for latency-critical flows.

Cloud-native serverless architectures with millisecond-level function invocations generate fundamentally different traffic patterns. Our temporal features assume flow durations of seconds to minutes, making them inappropriate for functions completing in 10 to 100 milliseconds. Necessary modifications include re-engineering temporal features for microsecond granularity, developing invocation-level rather than flow-level analysis, integrating with serverless platform telemetry, and designing lightweight models suitable for edge-based inference.

\subsection{I.3 Zero-Day Attack Detection Limitations}

Our 72 to 76 percent detection rate for zero-day attacks reflects fundamental limitations of supervised learning approaches. For signature-based approaches, theoretical detection probability for zero-day attacks is bounded by expression depending on model capacity and feature overlap. Under reasonable assumptions of 30 to 40 percent feature overlap, this yields detection probability in range 60 to 65 percent. FedGTD achieves 72 to 76 percent, exceeding theoretical bound by leveraging cross-domain knowledge transfer.

Improvement strategies include hybrid detection architecture combining supervised FedGTD for known attack classification with unsupervised anomaly detection for zero-day identification, using ensemble predictions with weighted voting based on confidence scores. Few-shot meta-learning trains meta-learners that adapt to novel attacks with minimal examples, implementing Model-Agnostic Meta-Learning for rapid adaptation, and using prototypical networks to learn attack family embeddings.

Continual learning implements online learning to incorporate zero-day examples immediately upon discovery, uses experience replay to prevent catastrophic forgetting of known attacks, and applies elastic weight consolidation to preserve knowledge while adapting to new threats.

\section{Appendix J: Extended Future Research Directions}

\subsection{J.1 Graph Neural Networks for Attack Propagation}

Current approaches treat security domains independently. Graph neural networks offer powerful tools for capturing cross-domain attack propagation through heterogeneous graph construction where nodes represent organizations, domains, security assets, and detected threats, while edges capture network connections, data flows, trust relationships, and temporal correlations. Node features would encode security posture, reputation scores, and attack history, while edge features represent connection strength, latency, and trust levels.

GNN-based attack prediction could leverage Graph Attention Networks to learn relative importance of different connections through learned attention weights, implement temporal graph convolutions to model dynamic attack evolution over time, apply heterogeneous GNN architectures to handle different node and edge types appropriately, and use link prediction techniques to forecast likely attack propagation paths enabling proactive defense.

This approach would enable predicting attack lateral movement before it occurs allowing proactive defensive measures, identifying critical vulnerabilities in organizational network topology through centrality analysis, optimizing defensive resource allocation based on quantified attack propagation risk, and enabling proactive threat hunting across federated organizations.

\subsection{J.2 Automated Feature Discovery via Transformers}

Manual feature engineering for each security domain is labor-intensive and may miss important patterns. Transformer architectures offer promising solutions through self-attention mechanisms applied to network flows. Treating network flow packets as sequences and applying multi-head self-attention to learn packet-level relationships, using positional encodings to capture temporal ordering information, and generating learned embeddings as features for downstream classification would eliminate significant manual effort.

Cross-domain attention mechanisms represent another valuable direction. Implementing cross-attention between security domains to discover shared patterns, learning domain-invariant representations through adversarial training approaches, and using attention weights as interpretability mechanism would enhance both performance and understanding.

These transformer-based approaches would eliminate manual feature engineering overhead currently required when adapting to new security domains, discover non-obvious cross-domain correlations that human experts might miss, improve zero-day attack detection through learned abstractions that generalize beyond specific attack signatures, and enable rapid adaptation to new security domains with minimal domain expertise.

\subsection{J.3 Continual Meta-Learning for Zero-Day Adaptation}

Extending FedGTD with meta-learning capabilities would enable rapid adaptation to novel attacks with minimal labeled examples. Model-Agnostic Meta-Learning integration offers promising approach through meta-training on diverse attack families to learn good parameter initializations enabling rapid fine-tuning, fine-tuning on few examples of novel attacks with small gradient updates leveraging the meta-learned initialization, using task-specific learning rates optimized during meta-training, and maintaining separate meta-parameters for each security domain.

Prototypical networks provide alternative approach for attack family classification. Learning embedding space where attacks cluster by family through metric learning, computing prototype vectors for each known attack family as centroids, classifying novel attacks by distance to nearest prototype enabling zero-shot or few-shot classification, and updating prototypes incrementally as new attacks are discovered would create continuously adapting detection system.

These meta-learning approaches could improve zero-day detection rates from current 72 to 76 percent up to 85 to 90 percent through better generalization, reduce time-to-detection for novel attacks from days to hours enabling rapid response, enable federated few-shot learning across organizations where one organization's discovery rapidly propagates to all participants, and maintain robust performance under concept drift and adversarial evolution.

\subsection{J.4 Explainable AI for Security Operations}

Current models lack interpretability. Future work should integrate explainability mechanisms. Attention-based explanations can visualize attention weights showing which features influenced predictions enabling analysts to understand model reasoning, generate natural language explanations providing context without requiring technical expertise, and provide counterfactual explanations indicating conditions under which prediction would change.

SHAP values for federated models offer another promising direction. Computing Shapley values to attribute prediction contributions to individual features quantifying each feature's impact, aggregating SHAP explanations across federated clients while preserving privacy, and generating global explanations showing important features per attack type revealing general patterns would significantly improve model transparency.

Example-based explanations complement these approaches by retrieving similar historical attacks from secure distributed databases enabling case-based reasoning, showing analysts how current alerts relate to known attack campaigns providing valuable context, and enabling interactive refinement of detection rules based on analyst feedback.

These explainability mechanisms would increase analyst trust and adoption of AI-based detection systems, accelerate incident investigation through contextual explanations, enable regulatory compliance requiring explainable AI decisions, and facilitate knowledge transfer between junior and senior security analysts.

\subsection{J.5 Quantum-Resistant Cryptography}

As quantum computers advance, current cryptographic primitives securing federated aggregation will become vulnerable. Primary threats include RSA-based client authentication broken by Shor's algorithm, elliptic curve cryptography remaining quantum vulnerable, and hash-based commitments partially vulnerable to Grover's algorithm.

Lattice-based cryptography offers viable solutions with proposals to replace RSA authentication using NIST-approved CRYSTALS-Dilithium signatures and employ Kyber for key encapsulation in secure aggregation. These replacements introduce expected performance overhead of 20 to 30 percent increase in computation time. Hash-based signatures, particularly SPHINCS+ for one-time client commitments and hash chains for reputation score verification impose minimal overhead.

The transition timeline should follow phased approach. Near-term deployment from 2025 to 2027 should focus on hybrid classical and post-quantum implementations allowing gradual migration. Medium-term plans from 2027 to 2030 should complete full migration to post-quantum primitives across all system components. Long-term development beyond 2030 should explore integration with quantum key distribution networks.

\subsection{J.6 Privacy-Preserving Threat Intelligence Sharing}

Extending FedGTD to enable secure sharing of threat intelligence beyond model parameters represents important future direction. Homomorphic encryption schemes can encrypt attack signatures using fully homomorphic encryption, enabling similarity matching on encrypted signatures while sharing positive matches without revealing full signature details. Secure multi-party computation protocols can compute aggregate attack statistics including prevalence and severity without revealing individual organizational data.

Differential privacy mechanisms applied to threat intelligence before sharing with partners can balance privacy concerns preventing attribution to specific organizations with utility requirement of providing actionable intelligence. Local differential privacy provides additional protection for highly sensitive threat indicators.

These mechanisms would enable industry-wide threat intelligence sharing despite competitive concerns, accelerating collective response to emerging threats, reducing time-to-detection for distributed attack campaigns, and ensuring compliance with information sharing regulations such as CISA directives and NIS2 Directive.

\rev{
\section{Appendix K: Extended Game-Theoretic Sensitivity Analysis}
\label{sec:extended_sensitivity}

\subsection{K.1 One-at-a-Time Parametric Analysis}

This appendix provides extended details on the sensitivity analysis summarized in Table~10 of the main paper. For each of the five game-theoretic parameters (detection reward coefficient $\alpha^{(d)}$, false-positive cost coefficient $\beta^{(d)}$, cross-domain synergy weight $\gamma$, discount factor $\gamma_{\text{disc}}$, and imbalance ratio $\rho_d$), we varied the parameter across ten equally spaced values spanning $\pm 50\%$ of its baseline while holding all other parameters fixed at their baseline values, and computed the resulting Nash equilibrium by running the full FedGTD algorithm for 200 rounds.

The detection reward coefficient $\alpha^{(d)}$ controls how strongly defenders are incentivized to detect attacks. As $\alpha^{(d)}$ increases, defenders adopt more aggressive detection thresholds, improving true positive rates at the cost of higher false positives. The equilibrium detection accuracy increases monotonically from 94.1\% at $\alpha^{(d)} = 0.5 \alpha_0$ to 96.8\% at $\alpha^{(d)} = 1.5 \alpha_0$, where $\alpha_0$ is the baseline value. The false-positive rate increases from 2.1\% to 5.3\% across the same range, confirming the expected detection-versus-false-alarm tradeoff.

The false-positive cost coefficient $\beta^{(d)}$ has the largest influence on equilibrium outcomes. Increasing $\beta^{(d)}$ from 0.5 to 1.5 times its baseline causes the equilibrium detection accuracy to decrease from 97.8\% to 93.2\%, a swing of 4.6 percentage points. This parameter effectively sets the operating point on the ROC curve: organizations with low tolerance for false alarms should set $\beta^{(d)}$ high, accepting some reduction in detection rates.

The remaining parameters have comparatively modest influence. The cross-domain synergy weight $\gamma$ affects equilibrium accuracy by at most 0.8\%, confirming that the adversary's cross-domain coordination is not the primary driver of detection performance. The discount factor $\gamma_{\text{disc}}$ has the smallest influence with a maximum deviation of 0.4\%, consistent with the observation that our federated learning operates over finite horizons where discounting plays a limited role.

\subsection{K.2 Monte Carlo Joint Sensitivity Analysis}

To assess robustness under simultaneous parameter perturbation, we sampled 1000 parameter vectors uniformly from the hypercube defined by $\pm 25\%$ of each baseline parameter. For each sample, we computed the Nash equilibrium and recorded the overall detection accuracy. The resulting distribution has mean 95.8\% with standard deviation 1.2\%, and the 5th percentile is 93.7\%. No parameter combination within the tested range produced accuracy below 92.1\%, indicating that the framework maintains strong performance across a wide range of operational conditions.

We computed partial rank correlation coefficients (PRCC) to identify the most influential parameters under joint variation. The PRCC analysis confirmed the findings from the one-at-a-time study: the false-positive cost coefficient has the largest absolute PRCC magnitude (0.72), followed by the detection reward (0.48), imbalance ratio (0.31), cross-domain synergy (0.18), and discount factor (0.09).

\section{Appendix E.4: Preprocessing Ablation Study}
\label{sec:preprocess_ablation}

To quantify the impact of individual preprocessing steps on detection performance, we systematically removed each preprocessing component and measured the resulting degradation across all three domains.

\begin{table}[H]
\centering
\caption{Preprocessing Ablation: Detection Accuracy (\%) After Removing Individual Steps}
\label{tab:preprocess_ablation}
\small
\begin{tabular}{lccc}
\toprule
\textbf{Removed Step} & \textbf{Edge-IIoT} & \textbf{Container} & \textbf{SOC} \\
\midrule
None (full pipeline) & 95.7 & 96.3 & 96.9 \\
Protocol-specific encoding & 91.2 ($-4.5$) & 96.1 ($-0.2$) & 96.8 ($-0.1$) \\
Temporal feature extraction & 94.3 ($-1.4$) & 95.4 ($-0.9$) & 93.1 ($-3.8$) \\
Class-weight calibration & 93.8 ($-1.9$) & 92.7 ($-3.6$) & 88.4 ($-8.5$) \\
Feature normalization & 94.1 ($-1.6$) & 94.8 ($-1.5$) & 95.2 ($-1.7$) \\
Missing value imputation & 95.5 ($-0.2$) & 95.9 ($-0.4$) & 94.6 ($-2.3$) \\
Entity embedding (SOC only) & 95.7 ($\pm 0.0$) & 96.3 ($\pm 0.0$) & 91.4 ($-5.5$) \\
\bottomrule
\end{tabular}
\end{table}

Table~\ref{tab:preprocess_ablation} reveals that the most critical preprocessing step varies by domain. For Edge-IIoT, protocol-specific encoding is essential, as removing it degrades accuracy by 4.5\%. For the Container domain, class-weight calibration is most important ($-3.6\%$ degradation), reflecting the severe 15.7:1 imbalance. For the SOC domain, class-weight calibration produces the largest degradation ($-8.5\%$), consistent with the extreme 99:1 imbalance; entity embedding is also critical ($-5.5\%$), as it captures the semantic relationships among the 33 entity types that are central to incident triage. Feature normalization has a consistent moderate impact across all domains, confirming its role as a general-purpose preprocessing step. These results validate our domain-specific preprocessing pipeline design and quantify the contribution of each component to overall detection performance.
}

\section{Conclusion}

This supplementary material provides comprehensive details supporting the main paper's contributions. We have presented extended related work demonstrating how FedGTD addresses gaps left by prior approaches. The detailed privacy analysis establishes that our framework meets stringent requirements including GDPR, HIPAA, and SOX compliance while maintaining detection effectiveness.

Complete mathematical proofs for all four main theorems establish rigorous theoretical foundations. Theorem 1 proves existence and uniqueness of solutions to coupled stochastic differential equations. Theorem 2 demonstrates existence of Nash equilibrium despite extreme heterogeneity and imbalance. Theorem 3 provides simultaneous privacy and robustness guarantees with explicit bounds. Theorem 4 establishes almost-sure convergence with concrete convergence rates under domain-adaptive learning. These proofs demonstrate that FedGTD rests on solid mathematical foundations.

Extended adversarial training procedures tailored to each security domain show how protocol-aware perturbations for Edge-IIoT, CVE-targeted training for containers, and entity-aware generation for SOC systems contribute to robust detection. The detailed experimental configurations ensure reproducibility through complete specification. Fair baseline comparison protocols eliminate confounding factors.

Additional experimental results including precision-recall curves, detailed confusion matrices, and comprehensive statistical testing validate significance of improvements. The extended scalability analysis identifies quadratic Byzantine detection as primary bottleneck while proposing hierarchical aggregation, sampling-based detection, and approximate similarity via locality-sensitive hashing as viable solutions.

Detailed deployment considerations provide practical guidance. The cost-benefit analysis demonstrates positive return on investment within 8 to 12 months. Integration guidelines provide step-by-step procedures. Operational best practices cover model update schedules, Byzantine client handling procedures, and privacy budget management strategies.

The comprehensive limitations analysis honestly assesses current boundaries through detailed examination of assumption dependencies and failure modes. Understanding when assumptions fail and what happens when Byzantine fractions exceed theoretical thresholds enables practitioners to anticipate and mitigate potential issues. The analysis of generalization limitations identifies specific adaptations required for deployment in various environments.

Future research directions spanning graph neural networks, transformer-based automated feature discovery, continual meta-learning, explainable AI, quantum-resistant cryptography, and privacy-preserving threat intelligence sharing chart ambitious but achievable research trajectories. These directions address current limitations while opening new capabilities that would further enhance federated security.

The material in this appendix demonstrates that FedGTD represents a mature, production-ready framework with strong theoretical foundations, comprehensive empirical validation, and clear pathways for continued development. We hope this supplementary material serves as a valuable resource for researchers seeking to understand the theoretical underpinnings, practitioners aiming to deploy federated security systems, and future investigators building upon this foundation to advance multi-cloud security.

\bibliographystyle{plain}
\begin{thebibliography}{99}

\bibitem{alpcan2011network}
Alpcan, T., Başar, T.: Network Security: A Decision and Game-Theoretic Approach. Cambridge University Press (2011)

\bibitem{roy2010survey}
Roy, S., Ellis, C., Shiva, S., Dasgupta, D., Shandilya, V., Wu, Q.: A survey of game theory as applied to network security. In: Proc. 43rd Hawaii Int'l Conf. System Sciences, pp. 1--10 (2010)

\bibitem{zhu2015game}
Zhu, Q., Başar, T.: Game-theoretic methods for robustness, security, and resilience of cyberphysical control systems. IEEE Control Systems Magazine \textbf{35}(1), 46--65 (2015)

\bibitem{preuveneers2018federated}
Preuveneers, D., Rimmer, V., Tsingenopoulos, I., Spooren, J., Joosen, W., Ilie-Zudor, E.: Chained anomaly detection models for federated learning: An intrusion detection case study. Applied Sciences \textbf{8}(12), 2663 (2018)

\bibitem{alhawawreh2018federated}
Al-Hawawreh, M., Sitnikova, E., Aboutorab, N.: X-IIoTID: A connectivity-agnostic and device-agnostic intrusion data set for industrial internet of things. IEEE Internet of Things Journal \textbf{9}(5), 3962--3977 (2022)

\bibitem{wang2024cloudfl}
Wang, L., Zhang, X., Liu, M.: CloudFL: Federated learning for multi-cloud security intelligence. IEEE Transactions on Cloud Computing \textbf{12}(3), 1245--1259 (2024)

\bibitem{chen2024adaptive}
Chen, Y., Li, H., Park, J.: Adaptive federated learning for container security in heterogeneous cloud environments. Computer Networks \textbf{238}, 110087 (2024)

\bibitem{liu2024hierarchical}
Liu, J., Zhao, S., Wang, K.: Hierarchical federated learning for IoT security with domain adaptation. IEEE Internet of Things Journal \textbf{11}(8), 13245--13258 (2024)

\bibitem{blanchard2017machine}
Blanchard, P., El Mhamdi, E.M., Guerraoui, R., Stainer, J.: Machine learning with adversaries: Byzantine tolerant gradient descent. In: Advances in Neural Information Processing Systems, pp. 119--129 (2017)

\bibitem{yin2018byzantine}
Yin, D., Chen, Y., Kannan, R., Bartlett, P.: Byzantine-robust distributed learning: Towards optimal statistical rates. In: Int'l Conf. Machine Learning, pp. 5650--5659 (2018)

\bibitem{zhang2024balance}
Zhang, C., Li, Y., Liu, X.: BALANCE: Byzantine-resilient decentralized federated learning using local model as reference for similarity. In: Proc. 2024 ACM SIGSAC Conf. Computer and Communications Security, pp. 1876--1891 (2024)

\bibitem{zhou2024robust}
Zhou, X., Patel, R., Chen, M.: RobustFL: Byzantine-resilient federated learning with improved aggregation mechanisms. IEEE Transactions on Information Forensics and Security \textbf{19}, 4567--4580 (2024)

\bibitem{patel2024secure}
Patel, R., Singh, M., Zhang, L.: SecureAgg+: Enhanced Byzantine-robust aggregation for federated learning. In: Proc. IEEE INFOCOM, pp. 1234--1242 (2024)

\bibitem{singh2024byzantine}
Singh, M., Chen, A., Kumar, P.: Byzantine-resilient federated learning for edge computing environments. IEEE Transactions on Mobile Computing \textbf{23}(7), 3456--3469 (2024)

\bibitem{mcmahan2017communication}
McMahan, B., Moore, E., Ramage, D., Hampson, S., y Arcas, B.A.: Communication-efficient learning of deep networks from decentralized data. In: Proc. 20th Int'l Conf. Artificial Intelligence and Statistics, pp. 1273--1282 (2017)

\bibitem{li2020federated}
Li, T., Sahu, A.K., Talwalkar, A., Smith, V.: Federated learning: Challenges, methods, and future directions. IEEE Signal Processing Magazine \textbf{37}(3), 50--60 (2020)

\bibitem{karimireddy2020scaffold}
Karimireddy, S.P., Kale, S., Mohri, M., Reddi, S.J., Stich, S.U., Suresh, A.T.: SCAFFOLD: Stochastic controlled averaging for federated learning. In: Int'l Conf. Machine Learning, pp. 5132--5143 (2020)

\bibitem{geyer2017differentially}
Geyer, R.C., Klein, T., Nabi, M.: Differentially private federated learning: A client level perspective. arXiv preprint arXiv:1712.07557 (2017)

\bibitem{zhang2024privacy}
Zhang, H., Liu, K., Wang, J.: Privacy-preserving federated learning with adaptive noise injection. IEEE Transactions on Dependable and Secure Computing \textbf{21}(4), 2134--2147 (2024)

\bibitem{kumar2024dp}
Kumar, A., Patel, S., Chen, M.: Domain-specific differential privacy for heterogeneous federated learning. In: Proc. ACM CCS, pp. 567--579 (2024)

\bibitem{bottou2018sgd}
Bottou, L., Curtis, F.E., Nocedal, J.: Optimization methods for large-scale machine learning. SIAM Review \textbf{60}(2), 223--311 (2018)

\bibitem{manshaei2013game}
Manshaei, M.H., Zhu, Q., Alpcan, T., Başar, T., Hubaux, J.-P.: Game theory meets network security and privacy. ACM Computing Surveys \textbf{45}(3), 1--39 (2013)

\bibitem{fang2020local}
Fang, M., Cao, X., Jia, J., Gong, N.: Local model poisoning attacks to Byzantine-robust federated learning. In: 29th USENIX Security Symposium, pp. 1605--1622 (2020)

\bibitem{aghassi2006robust}
\rev{Aghassi, M., Bertsimas, D.: Robust game theory. Mathematical Programming, Series B \textbf{107}(1--2), 231--273 (2006)}

\bibitem{saltelli2004sensitivity}
\rev{Saltelli, A., Tarantola, S., Campolongo, F., Ratto, M.: Sensitivity Analysis in Practice: A Guide to Assessing Scientific Models. Wiley (2004)}

\end{thebibliography}

\end{document}